% Sección 9 de la guía: conclusiones. Rol 4.
\section{Conclusiones}
\label{sec:conclusiones}
Se predijo la concentración horaria de PM$_{2{,}5}$ una hora adelante en la estación Santa Anita con una partición cronológica y control explícito de fuga de información. Las respuestas a las preguntas de investigación son las siguientes.

\textbf{PI1.} Los modelos de árboles reducen el MAE de la persistencia en 9,5--11,6~\% en validación y alrededor de 14~\% en prueba (2,41 y 2,43 frente a 2,83~\textmu g/m\textsuperscript{3}), y Ridge en 7~\%. Los intervalos de la diferencia excluyen el cero, pero la ganancia es pequeña en términos absolutos.

\textbf{PI2.} El calendario aporta solo a los árboles (2 a 3~\% en validación) y por medio de la hora del día; el día de la semana no aporta y en los modelos lineales el calendario empeora ligeramente el MAE. La hipótesis se confirma de forma parcial.

\textbf{PI3.} Sí se degrada: en el decil superior el MAE es de dos a cuatro veces el del resto y por encima de 50~\textmu g/m\textsuperscript{3} todos los modelos subestiman. Los árboles conservan su ventaja relativa sobre la persistencia, pero ninguno resuelve la subestimación de los picos.

En conjunto, el historial propio de la serie permite mejoras consistentes pero modestas a una hora, y no basta para anticipar los episodios de mayor relevancia sanitaria. Como trabajo futuro se propone incorporar variables meteorológicas, evaluar otras estaciones y horizontes más largos, probar modelos que atiendan la cola superior (pérdidas ponderadas o regresión cuantílica), explorar estrategias de imputación que no usen información futura y repetir la evaluación con más periodos para confirmar la estabilidad de los resultados.
